\documentclass[10pt,a4paper]{amsart}
\usepackage[margin=19mm]{geometry}
\usepackage{amsmath,amssymb,mathtools}
\usepackage[scr=rsfs,cal=euler]{mathalfa}
\usepackage{xfrac}
\usepackage[unicode,hidelinks,bookmarks=false]{hyperref}
\newcommand{\ie}{i.e.\ }
\newcommand{\cf}{cf.\ }
\newcommand{\resp}{respectively\ }
\newcommand{\Subsection}{\subsection}
\providecommand{\nobreakdash}{\nobreak\hbox{-}}
\setlength{\emergencystretch}{2em}
\multlinegap0pt
\usepackage[shortlabels]{enumitem}
\usepackage{amssymb}
\usepackage{mathtools}
\newtheorem{theorem}{Theorem}
\newtheorem{lemma}{Lemma}
\newtheorem{corollary}{Corollary}
\newtheorem{proposition}{Proposition}
\newcommand{\Acal}{\mathcal{A}}
\newcommand{\N}{\mathbb{N}}
\newcommand{\R}{\mathbb{R}}
\newcommand{\Rxd}{\R[x_1,\dots,x_d]}
\title{The $K$-Moment Problem :\\ a detailed introduction}
\author{Malik Amir}
\date{Source: arXiv:2604.12265v1 (CC BY 4.0)}
\begin{document}
\maketitle

\noindent\emph{Source and licence.} The $K$-Moment Problem :\\ a detailed introduction. Version arXiv:2604.12265v1 (CC BY 4.0), distributed by arXiv under the Creative Commons Attribution 4.0 International licence, http://creativecommons.org/licenses/by/4.0/, as recorded in the arXiv licence metadata for this exact version and shown on the abstract page https://arxiv.org/abs/2604.12265v1. Full licence notice: \texttt{licenses/arxiv-2604.12265-CC-BY-4.0.txt}.\par\smallskip

\noindent\textit{Editorial scope.} Counted unit: the article's proposition prop:compactimpliesarch (source0.tex lines 369--371). The article's complete original proof of it is delivered with it (source0.tex lines 373--407, one \texttt{proof} environment of 35 lines). No mathematics is written, repaired or re-derived: the statement and the proof are the article's own lines, in the article's own notation, and no retained line is substituted or re-flowed.  Retained uncounted as the source-local context the proof needs: lines 250--267; lines 289--295; lines 320--333.  Nothing was omitted from any retained span, so the package carries no omission record.  No retained line contains a citation, so the wrapper carries no bibliography and no bibliography entry.  No retained line contains a figure environment, an image-inclusion command or drawing code, so no image is delivered and the package declares no asset dependency.  Editorial formatting only, in addition: an AMS \texttt{amsart} preamble that restates the article's own declarations and macros used by the retained lines, namely the environments \texttt{theorem}, \texttt{lemma}, \texttt{corollary}, \texttt{proposition} on separate counters, together with the standard packages those lines need.  The retained environments print in this document as Theorem 1, Lemma 1, Corollary 1, Proposition 1 in order of appearance, whereas the article prints the same items with the numbering of its own full document.  The complete native source of the article as distributed in the arXiv e-print for this exact version is bundled unchanged as \texttt{sources/198394/source0.tex}.
\begin{theorem}[Krivine--Stengle]\label{thm:ks}
Let \(f=(f_1,\dots,f_s)\subset \Rxd\), and let \(g\in\Rxd\). Then the following statements hold.
\begin{enumerate}[label=\textup{(\arabic*)}, leftmargin=2em]
\item \(g>0\) on \(K(f)\) if and only if there exist \(p,q\in T(f)\) such that
\[
pg=1+q.
\]
\item \(g\ge 0\) on \(K(f)\) if and only if there exist \(p,q\in T(f)\) and \(m\in\N\) such that
\[
pg=g^{2m}+q.
\]
\item \(g=0\) on \(K(f)\) if and only if there exists \(m\in\N\) such that
\[
-g^{2m}\in T(f).
\]
\item \(K(f)=\varnothing\) if and only if \(-1\in T(f)\).
\end{enumerate}
\end{theorem}

\begin{lemma}\label{lem:boundedcriterion}
Let \(Q\) be a quadratic module in \(\Acal\) and let \(a\in\Acal\). Then the following are equivalent.
\begin{enumerate}[label=\textup{(\roman*)}, leftmargin=2em]
\item \(a\in \Acal_b(Q)\).
\item There exists \(\lambda>0\) such that \(\lambda^2-a^2\in Q\).
\end{enumerate}
\end{lemma}

\begin{corollary}\label{cor:archpoly}
For a quadratic module \(Q\subseteq\Rxd\), the following are equivalent.
\begin{enumerate}[label=\textup{(\roman*)}, leftmargin=2em]
\item \(Q\) is Archimedean.
\item There exists \(\lambda>0\) such that
\[
\lambda-\sum_{j=1}^d x_j^2\in Q.
\]
\item For each \(j=1,\dots,d\), there exists \(\lambda_j>0\) such that
\[
\lambda_j-x_j^2\in Q.
\]
\end{enumerate}
\end{corollary}

\begin{proposition}\label{prop:compactimpliesarch}
If \(K(f)\) is compact, then \(T(f)\) is Archimedean.
\end{proposition}

\begin{proof}
Let
\[
g(x):=(1+x_1^2)\cdots(1+x_d^2).
\]
Since \(K(f)\) is compact, \(g\) is bounded on \(K(f)\). Choose \(\lambda>0\) such that
\[
\lambda^2>g(x)^2 \qquad \text{for all } x\in K(f).
\]
By Theorem \ref{thm:ks}\textup{(1)}, there exist \(p,q\in T(f)\) such that
\[
p(\lambda^2-g^2)=1+q.
\]
Fix \(n\in\N_0\). Multiplying by \(g^{2n}\) and using that \(T(f)\) is a quadratic module, we obtain
\[
g^{2n+2}p=\lambda^2g^{2n}p-g^{2n}(1+q)\preceq \lambda^2g^{2n}p,
\]
where \(a\preceq b\) means \(b-a\in T(f)\). By induction,
\[
g^{2n}p\preceq \lambda^{2n}p.
\]
Since \(g^{2n}(q+pg^2)\in T(f)\), we also have
\[
g^{2n}\preceq g^{2n}(1+q+pg^2)=\lambda^2g^{2n}p\preceq \lambda^{2n+2}p.
\]
Because \(p\) contains only finitely many monomials, there exist \(c>0\) and \(k\in\N\) such that \(p\preceq c g^k\). Taking \(n=k\) yields
\[
g^{2k}\preceq c\lambda^{2k+2}g^k.
\]
Hence
\[
\left(g^k-\tfrac12 c\lambda^{2k+2}\right)^2\preceq \left(\tfrac12 c\lambda^{2k+2}\right)^2,
\]
so Lemma \ref{lem:boundedcriterion} shows that \(g^k\) is \(T(f)\)-bounded. Since \(\pm x_j\preceq g^k\) for each \(j\), every coordinate function is \(T(f)\)-bounded, and Corollary \ref{cor:archpoly} implies that \(T(f)\) is Archimedean.
\end{proof}

\end{document}
